% ID: SEL-D0084-OA-E04
% ESTADO: cotejada. Segunda pasada visual y borrador compilado revisados.
% Fuente DOC: Selectividad/sel_2012_matematicas_aplicadas/6 MATEMÁT. APLIC. CC. SS/EXAMEN 6-ANDALUCÍA.doc
% SHA-256 original: 61f8d5d809881c49e9e9212281f7028516338e026524e55d572ab691ec5f416d
% Página convertida: 1; opción A; ejercicio 4.

\textbf{(2.5 puntos)} Un índice para calibrar la madurez lectora de los alumnos de primaria se distribuye según una ley Normal con desviación típica 2. Elegida una muestra de 18 alumnos en un centro de primaria, se obtiene una media muestral de $10.8$ en dicho índice. Mediante el uso de un contraste de hipótesis, \textquestiondown{}se puede aceptar, con un nivel de significación del $1\%$, la hipótesis nula de que la media del índice de madurez lectora de los alumnos de este centro no es inferior a 11?
